% Resumen — máx. 1 PÁGINA [REGLAS_FORMATO §3, RUBRICA "Resumen / Abstract" 5 %].
% Debe contener: contexto, objetivos, metodología, resultados CON DATOS CUANTITATIVOS
% y conclusiones. Luego, entre 3 y 5 palabras o expresiones clave.
% Origen: v3.0 (2026-11-19), citas migradas a biblatex.
\chapter*{Resumen}
\addcontentsline{toc}{chapter}{Resumen}
\label{pag:resumen-ini}

% Estructura pedida por el profesor guía (16-11-2026): 4 párrafos.
% (1) introductorio · (2) objetivo general + algo de metodología · (3) resultados · (4) síntesis de conclusiones.

La Región del Maule enfrenta una creciente escasez hídrica que afecta a la producción agrícola, por lo que la reutilización de aguas residuales tratadas para riego surge como una fuente alternativa de agua.
En zonas rurales, los sistemas de tratamiento convencionales resultan complejos de implementar, mientras que las Soluciones Basadas en la Naturaleza, como los humedales construidos, ofrecen una alternativa descentralizada y de baja demanda energética.
En este contexto, el proyecto ANID/Anillos/ATE250033 contempla una estación experimental para tratar tres fuentes de agua residual de un predio frutícola de la Región del Maule: aguas servidas domésticas, aguas de lavado de agroquímicos y agua del estero Batuco.

El objetivo de este trabajo fue diseñar un sistema experimental de humedales construidos de flujo vertical para el tratamiento de esas tres fuentes y su futura reutilización en riego, integrando el sistema hidráulico y las necesidades estructurales para su implementación.
Para ello se caracterizaron las fuentes mediante seis campañas de muestreo en cinco puntos; se dimensionaron los humedales con el método del balance de oxígeno; se calcularon las redes por gravedad y las líneas a presión con la ecuación de Hazen-Williams; y se diseñaron los elementos de soporte con el método LRFD y el análisis sísmico de la NCh2369.

\pendiente{Párrafo de resultados con cifras: DQO y nitrógeno de cada fuente; caudal de la línea doméstica (\qty{0.75}{\cubic\meter\per\day}); \qty{6}{\square\meter} por línea en 6 IBC, con \qty{315}{\gram\per\day} de oxígeno frente a una demanda de \qty{111.8}{\gram\per\day}; ADT y bombas; perfiles de la plataforma y presión de contacto (cuando resuelva las preguntas al profesor).}

\pendiente{Párrafo de conclusiones en 3--4 oraciones: factibilidad del diseño, condiciones críticas (horizontalidad de la plataforma, impulsión del estero) y uso del expediente para construir la estación.}

\noindent\textbf{Palabras clave:} Humedales Construidos, Soluciones Basadas en la Naturaleza, Diseño Hidráulico, Cálculo Estructural.
\label{pag:resumen-fin}
